\documentclass[12pt, a4paper]{article}
\usepackage[utf8]{inputenc}
\usepackage[T1]{fontenc}
\usepackage{lmodern}
\usepackage{amsmath}
\DeclareUnicodeCharacter{2212}{\ensuremath{-}}
\DeclareUnicodeCharacter{00D7}{\ensuremath{\times}}
\DeclareUnicodeCharacter{2192}{\ensuremath{\rightarrow}}
\DeclareUnicodeCharacter{0394}{\ensuremath{\Delta}}
\DeclareUnicodeCharacter{03B2}{\ensuremath{\beta}}
\DeclareUnicodeCharacter{03B7}{\ensuremath{\eta}}
\DeclareUnicodeCharacter{03B1}{\ensuremath{\alpha}}
\DeclareUnicodeCharacter{03B8}{\ensuremath{\theta}}
\DeclareUnicodeCharacter{03C4}{\ensuremath{\tau}}
\DeclareUnicodeCharacter{03C3}{\ensuremath{\sigma}}
\DeclareUnicodeCharacter{2264}{\ensuremath{\leq}}
\DeclareUnicodeCharacter{2208}{\ensuremath{\in}}
\DeclareUnicodeCharacter{2260}{\ensuremath{\neq}}
\DeclareUnicodeCharacter{00B1}{\ensuremath{\pm}}
\DeclareUnicodeCharacter{00B2}{\textsuperscript{2}}
\DeclareUnicodeCharacter{2081}{\textsubscript{1}}
\DeclareUnicodeCharacter{2082}{\textsubscript{2}}
\DeclareUnicodeCharacter{25B3}{\ensuremath{\triangle}}
\DeclareUnicodeCharacter{2032}{\ensuremath{\prime}}
\DeclareUnicodeCharacter{0304}{\={}}
\usepackage{booktabs}
\usepackage{longtable}
\usepackage{array}
\usepackage{calc}
\usepackage{float}
\usepackage{graphicx}
\PassOptionsToPackage{hyphens}{url}
\usepackage{hyperref}
\usepackage{xurl}
\usepackage{setspace}
\usepackage[margin=1in]{geometry}
\usepackage[font=it,justification=centering]{caption}
\usepackage{times}

\makeatletter
\setstretch{2.0}
\renewcommand{\arraystretch}{1.6}
\raggedbottom
\urlstyle{same}
\setlength{\emergencystretch}{3em}
\newenvironment{referenceshang}{%
  \par\begingroup
  \setlength{\parindent}{0pt}%
  \setlength{\parskip}{0.5\baselineskip}%
  \everypar{\hangindent=1.5em\hangafter=1\relax}%
  \ignorespaces
}{%
  \par\endgroup
}
\newcommand{\SITitleLine}[1]{\noindent {\Large\bfseries #1}\par}
\newcommand{\SIMajorHeading}[1]{\par\medskip\noindent {\large\bfseries #1}\par\smallskip}
\newcommand{\SIAnchorHeading}[2]{\clearpage\noindent \hypertarget{#1}{{\large\bfseries #2}}\par\smallskip\@afterindentfalse\@afterheading}
\makeatother

\begin{document}

% ============================================================================
% JPBM Round 2 Supplementary Information (2026-08-26).
% Carried over from the JBR Supportive Information (frozen baseline) with
% terminology-audit updates, plus two new sections: H (Supplementary Study 4,
% moved out of the main text) and I (additional methodological notes and
% sensitivity analyses receiving the detail consolidated out of the main text).
% Anonymized: no author-identifying information.
% ============================================================================

\SITitleLine{Supplementary Information}

\SITitleLine{Free or Cheap? Initial Price and Distribution Information,}

\SITitleLine{Failure Inference, and Parent-Brand Evaluation}

\SITitleLine{in Digital Luxury Brand Extensions}

\begin{quote}
\hyperlink{si:a1}{A-1. Study 1: Price manipulation texts}

\hyperlink{si:a2}{A-2. Study 2: Condition texts and stimulus images}

\hyperlink{si:a3}{A-3. Study 3a: Condition texts and displayed information}

\hyperlink{si:a4}{A-4. Study 3b: Condition texts and stimulus images}

\hyperlink{si:a5}{A-5. Study 3c: Condition texts}

\hyperlink{si:a6}{A-6. Supplementary Study 4: Condition texts}

\hyperlink{si:b}{B. Study 1: ANCOVA with covariates}

\hyperlink{si:c}{C. Single-paper meta-analysis across Studies 3a--3c}

\hyperlink{si:d}{D. Study 3a: Exploratory serial mediation analysis}

\hyperlink{si:e}{E. Preliminary Field Study: OpenSea floor-price analysis}

\hyperlink{si:f}{F. Preliminary Field Study: Brand classification}

\hyperlink{si:g}{G. Preliminary Field Study: X post analysis procedures and results}

\hyperlink{si:h}{H. Supplementary Study 4: Full report}

\hyperlink{si:i}{I. Additional methodological notes and sensitivity analyses}
\end{quote}

\SIAnchorHeading{si:a1}{A-1. Study 1: Price Manipulation (Four Levels)}

\noindent \textbf{High-price condition:} The luxury brand ``Gucci'' has launched a new NFT product. This NFT will be sold at a price higher than those of NFTs from various other brands.

\noindent \textbf{Comparable-price condition:} The luxury brand ``Gucci'' has launched a new NFT product. This NFT will be sold at a price comparable to those of NFTs from other luxury brands.

\noindent \textbf{Low-price condition:} The luxury brand ``Gucci'' has launched a new NFT product. This NFT will be sold at a price substantially lower than those of NFTs from various other brands.

\noindent \textbf{Free condition:} The luxury brand ``Gucci'' has launched a new NFT product. This NFT will be distributed for free.

\SIAnchorHeading{si:a2}{A-2. Study 2: Product Form (Physical vs. Digital) $\times$ Condition (Complimentary Magazine Gift vs. Comparable Price)}

\noindent \textbf{Physical $\times$ Complimentary magazine gift:} A Gucci mini-bag offered as a complimentary gift with a fashion magazine (limited to 1,000 units).

\noindent \textbf{Physical $\times$ Comparable price:} A Gucci mini-bag sold at regular retail at a price comparable to those of other luxury brands (limited to 1,000 units).

\noindent \textbf{Digital $\times$ Complimentary magazine gift:} A Gucci mini-bag NFT image offered as a complimentary gift with a fashion magazine (limited to 1,000 units).

\noindent \textbf{Digital $\times$ Comparable price:} A Gucci mini-bag NFT image sold at a price comparable to those of NFTs from other luxury brands (limited to 1,000 units).

\begin{figure}[H]
\centering
\includegraphics[width=0.49\linewidth]{figures/fig_study2_phy.png}\includegraphics[width=0.47\linewidth]{figures/fig_study2_virt.png}
\caption{Stimulus images for Study 2. Left: physical-product condition. Right: digital (NFT) condition.}
\end{figure}

\SIAnchorHeading{si:a3}{A-3. Study 3a: Condition Texts and Displayed Buyer and Market-Value Information}

\noindent The two initial conditions were identical to the digital conditions of Study 2 (Section A-2). The displayed information, framed as the pricing situation one month later, was:

\noindent \textbf{Purchased/comparable condition:} Pricing situation one month later / Gucci mini-bag NFT / Sale price: comparable to those of NFTs from other luxury brands / Order book (buyers): 0 (none) / △ At present, this NFT has almost no market value.

\noindent \textbf{Complimentary magazine-gift condition:} Pricing situation one month later / Gucci mini-bag NFT / Sale price: free / Order book (buyers): 0 (none) / △ At present, this NFT has almost no market value.

\begin{figure}[H]
\centering
\includegraphics[width=\linewidth]{figures/fig_study3a_demand.png}
\caption{Displayed-information screen used in Study 3a (complimentary condition shown). Participants viewed this screen after the launch phase.}
\end{figure}

\SIAnchorHeading{si:a4}{A-4. Study 3b: Image (Central GG Logo vs. Bee Image with a Small Logo) $\times$ Condition (Complimentary Fan-Book Gift vs. Comparable Price)}

\noindent \textbf{Launch phase (image $\times$ condition manipulation)}

\noindent \textbf{Central GG logo $\times$ Comparable price:} Gucci logo NFT. Limited to 1,000 units. Price comparable to those of NFTs from other luxury brands.

\noindent \textbf{Central GG logo $\times$ Complimentary fan-book gift:} Gucci logo NFT. Limited to 1,000 units. Free (complimentary with a fan book).

\noindent \textbf{Bee image $\times$ Comparable price:} Gucci bee-motif NFT. Limited to 1,000 units. Price comparable to those of NFTs from other luxury brands.

\noindent \textbf{Bee image $\times$ Complimentary fan-book gift:} Gucci bee-motif NFT. Limited to 1,000 units. Free (complimentary with a fan book).

\noindent \textbf{Displayed buyer and market-value information:} (All conditions) Order book (buyers): 0 (none) / △ At present, this NFT has almost no market value.

\begin{figure}[H]
\centering
\includegraphics[width=0.39\linewidth]{figures/fig_study3b_fl_HIGH.png}\includegraphics[width=0.56\linewidth]{figures/fig_study3b_fl_low.png}
\caption{Stimulus images used in Study 3b. Left: image dominated by a central GG logo. Right: bee image with a small logo.}
\end{figure}

\SIAnchorHeading{si:a5}{A-5. Study 3c: Source (Official Brand Website vs. News Article) $\times$ Condition (Free vs. Comparable Price)}

\noindent \textbf{Launch phase (source $\times$ condition manipulation; all material text-only)}

\noindent \textbf{Official brand website $\times$ Comparable price:} louisvuitton.com / LOUIS VUITTON / OFFICIAL / January 15, 2025. Announcing the digital collection ``LOUIS VUITTON NFT.'' Limited to 1,000 units. Price comparable to those of NFTs from other luxury brands.

\noindent \textbf{Official brand website $\times$ Free:} louisvuitton.com / LOUIS VUITTON / OFFICIAL / January 15, 2025. Announcing the digital collection ``LOUIS VUITTON NFT.'' Limited to 1,000 units. Distributed for free.

\noindent \textbf{News article $\times$ Comparable price:} Fashion Tech News (NFT, January 15, 2025 \textbar{} Editorial). Louis Vuitton to enter the NFT market. Limited to 1,000 units. Price comparable to those of NFTs from other luxury brands.

\noindent \textbf{News article $\times$ Free:} Fashion Tech News (NFT, January 15, 2025 \textbar{} Editorial). Louis Vuitton to enter the NFT market. Limited to 1,000 units. Distributed for free.

\noindent \textbf{Displayed buyer and market-value information (market situation one month later):} All four conditions displayed: Initial price (comparable/free) / Buyers: 0 (none) / △ There are currently no buyers.

\SIAnchorHeading{si:a6}{A-6. Supplementary Study 4: Condition Texts (Favorable Displayed Information)}

\noindent \textbf{Displayed information (pricing situation one month later)}

\noindent \textbf{Comparable-price condition:} Pricing situation one month later / Gucci bee-motif NFT / Sale price: comparable to those of NFTs from other luxury brands / Order book (buyers for this NFT): 1,343 (very high) / △ Demand for this Gucci NFT is very high and its price is rising sharply.

\noindent \textbf{Free condition:} Pricing situation one month later / Gucci bee-motif NFT / Sale price: free / Order book (buyers for this NFT): 1,343 (very high) / △ Demand for this Gucci NFT is very high and its price is rising sharply.

\emph{Note} (exclusion criterion). Among the exclusions reported in Section H, the removal of respondents who scored at the scale ceiling (10) at both \emph{T}0 and \emph{T}1 was a criterion specified in the preregistration (AsPredicted \#262807), intended to exclude cases in which change could not be assessed because of a ceiling effect.

\emph{Note.} Brand names and logos appearing in this supplementary material are trademarks of the respective companies and are used solely for academic research purposes.

\SIAnchorHeading{si:b}{B. Study 1: ANCOVA with Covariates}

To assess the robustness of the ANOVA results in Study 1, we conducted ANCOVAs controlling for individual-difference variables. The covariates were self--brand connection (SBC; 6-item mean), NFT purchase experience, and Gucci purchase experience. The analytical sample was the same 670 participants reported in the main text (after excluding four who failed the attention check).

\noindent \textbf{Loss of Brand Luxuriousness.} The main effect of price condition remained significant in the ANCOVA (\emph{F}(3, 663) = 36.23, \emph{p} \textless{} .001, partial $\eta^2$ = .14), essentially identical to the original ANOVA (\emph{F}(3, 666) = 35.17, \emph{p} \textless{} .001, partial $\eta^2$ = .14). Among the covariates, SBC was significant (\emph{F}(1, 663) = 13.58, \emph{p} \textless{} .001, partial $\eta^2$ = .02), and Gucci purchase experience was also significant (\emph{F}(1, 663) = 5.53, \emph{p} = .019, partial $\eta^2$ = .01). NFT purchase experience was not significant (\emph{F}(1, 663) = 0.29, \emph{p} = .591). Adjusted means were: high-price \emph{M} = 2.93 (\emph{SE} = 0.11), comparable \emph{M} = 3.67 (\emph{SE} = 0.11), low-price \emph{M} = 4.54 (\emph{SE} = 0.11), and free \emph{M} = 3.92 (\emph{SE} = 0.11), closely matching the unadjusted means (2.93, 3.69, 4.54, 3.91, respectively).

\noindent \textbf{Behavioral Intention.} The main effect of price condition was significant (\emph{F}(3, 663) = 93.90, \emph{p} \textless{} .001, partial $\eta^2$ = .30). Compared with the original ANOVA (\emph{F}(3, 666) = 70.30, \emph{p} \textless{} .001, partial $\eta^2$ = .24), the \emph{F} value increased after controlling for covariates, primarily because controlling for SBC reduced error variance. SBC had a strong positive effect on Behavioral Intention (\emph{F}(1, 663) = 134.16, \emph{p} \textless{} .001, partial $\eta^2$ = .17). NFT purchase experience (\emph{F}(1, 663) = 2.03, \emph{p} = .155) and Gucci purchase experience (\emph{F}(1, 663) = 2.54, \emph{p} = .111) were not significant. Adjusted means were: high-price \emph{M} = 2.84 (\emph{SE} = 0.10), comparable \emph{M} = 3.71 (\emph{SE} = 0.10), low-price \emph{M} = 4.08 (\emph{SE} = 0.10), and free \emph{M} = 5.15 (\emph{SE} = 0.10), closely matching the unadjusted means.

\noindent \textbf{Summary.} The effects of price condition reported in Study 1 are robust to controlling for self--brand connection with Gucci, NFT purchase experience, and Gucci purchase experience.

\SIAnchorHeading{si:c}{C. Single-Paper Meta-Analysis Across Studies 3a--3c}

To summarize the direction and moderation of the estimated indirect effects (condition $\rightarrow$ Failure Inference $\rightarrow$ Loss of Brand Luxuriousness) across Studies 3a through 3c, we conducted a single-paper meta-analysis (SPM) following the recommendations of McShane and B\"ockenholt (2025). SPM provides a statistical summary of evidence across multiple studies within a single paper, enabling effect-size estimation, quantification of between-study heterogeneity, and moderator testing.

We collected the estimated indirect effects from all conditions across Studies 3a through 3c (\emph{k} = 5; one effect from Study 3a, and two per study from Studies 3b and 3c, separated by presentation condition). The moderator was a presentation-condition coding: 0 for Study 3a's mini-bag stimulus, the bee image with a small logo in Study 3b, and the news article in Study 3c; 1 for the image with a central GG logo in Study 3b and the official brand website in Study 3c. This coding pools three different free or complimentary implementations, two brands, and two presentation manipulations (image and source); the pooled estimates should therefore be read as a summary of the direction and moderation of the study-level estimates, not as evidence of a single common effect. A random-effects meta-regression with this coding as the moderator was conducted; the heterogeneity parameter $\tau^2$ was estimated using the DerSimonian--Laird method (DerSimonian \& Laird, 1986).
% Note: this coding is numerically identical to the coding used at JBR submission
% (formerly labeled "flagshipness"); only the label has been changed to reflect
% the audited terminology. All estimates below are unchanged.

The overall random-effects model yielded a significant negative indirect-effect estimate (\ensuremath{\bar{\theta}} = -0.37, \emph{SE} = 0.17, 95\% CI {[}-0.69, -0.04{]}, \emph{z} = -2.22, \emph{p} = .026). Between-study heterogeneity was substantial (\emph{Q}(4) = 21.80, \emph{p} \textless{} .001, \emph{I}\textsuperscript{2} = 81.70\%, $\tau^2$ = 0.11).

The meta-regression revealed that the presentation-condition coding significantly accounted for heterogeneity in the indirect effect ($\beta_1$ = 0.53, \emph{SE} = 0.23, \emph{z} = 2.26, \emph{p} = .024). The moderator explained 64.90\% of the heterogeneity, and residual heterogeneity was substantially reduced (residual $\tau^2$ = 0.04). For conditions coded 0 (mini-bag, bee image, news article), the indirect effect was significantly negative (predicted value = -0.57, 95\% CI {[}-0.87, -0.28{]}). For conditions coded 1 (central GG logo, official brand website), the indirect effect was close to zero and not significant (predicted value = -0.04, 95\% CI {[}-0.39, 0.31{]}).

We also meta-analyzed the indices of moderated mediation from Studies 3b and 3c. The pooled index was significantly positive (\ensuremath{\overline{\mathrm{IMM}}} = 0.47, \emph{SE} = 0.16, 95\% CI {[}0.17, 0.78{]}, \emph{z} = 3.04, \emph{p} = .002), with negligible heterogeneity (\emph{Q}(1) = 0.03, \emph{p} = .869), indicating that the moderation pattern was consistent across the two studies despite their different manipulations.

\SIAnchorHeading{si:d}{D. Study 3a: Exploratory Serial Mediation Analysis Including the Single Monetization Item}

The main text theorizes that initial price and distribution information changes the interpretive context in which displayed buyer and market-value information is read. Study 3a administered a single exploratory item, ``I believe this NFT is a product intended to generate revenue for the company'' (7-point scale). This analysis was not preregistered and is exploratory; the item is a single author-developed item, not a validated construct measure.

As reported in the main text, the purchased/comparable condition (\emph{M} = 4.74, \emph{SD} = 1.32, \emph{n} = 72) scored higher on this item than the complimentary magazine-gift condition (\emph{M} = 3.86, \emph{SD} = 1.68, \emph{n} = 66; Welch's \emph{t}(123.25) = 3.40, \emph{p} = .001, \emph{d} = 0.59).

Building on this result, we estimated a serial mediation model: condition (X) $\rightarrow$ single monetization item (M1) $\rightarrow$ Failure Inference (M2) $\rightarrow$ Loss of Brand Luxuriousness (Y) (PROCESS Model 6; Hayes, 2022; 5,000 bootstrap resamples).

\noindent \textbf{Model 1: Predicting the single monetization item (M1).} The complimentary condition significantly reduced scores on the item ($a_1$ = -0.88, \emph{SE} = 0.26, \emph{t}(136) = -3.41, \emph{p} = .001, 95\% CI {[}-1.39, -0.38{]}; \emph{R}\textsuperscript{2} = .08).

\noindent \textbf{Model 2: Predicting Failure Inference (M2).} Both the direct effect of condition ($a_2$ = -0.93, \emph{SE} = 0.23, \emph{t}(135) = -4.05, \emph{p} \textless{} .001) and the effect of the monetization item ($d_{21}$ = 0.27, \emph{SE} = 0.07, \emph{t}(135) = 3.70, \emph{p} \textless{} .001) were significant (\emph{R}\textsuperscript{2} = .24).

\noindent \textbf{Model 3: Predicting Loss of Brand Luxuriousness (Y).} Failure Inference significantly increased Loss of Brand Luxuriousness ($b_2$ = 0.65, \emph{SE} = 0.09, \emph{t}(134) = 7.54, \emph{p} \textless{} .001). The effect of the monetization item was also significant ($b_1$ = -0.17, \emph{SE} = 0.08, \emph{t}(134) = -2.19, \emph{p} = .030). The direct effect of condition was not significant ($c'$ = 0.07, \emph{SE} = 0.24, \emph{t}(134) = 0.31, \emph{p} = .755; \emph{R}\textsuperscript{2} = .32).

Indirect-effect estimates are summarized in Table D1.
\setcounter{table}{0}
\renewcommand{\thetable}{D\arabic{table}}
\renewcommand{\theHtable}{D\arabic{table}}

\begin{table}[H]
\centering
\caption{Indirect Effects From the Exploratory Serial Mediation Analysis (Study 3a; 5,000 Bootstrap Resamples)}
\begin{tabular}{@{}
  >{\raggedright\arraybackslash}p{(\columnwidth - 6\tabcolsep) * \real{0.3846}}
  >{\raggedright\arraybackslash}p{(\columnwidth - 6\tabcolsep) * \real{0.1709}}
  >{\raggedright\arraybackslash}p{(\columnwidth - 6\tabcolsep) * \real{0.1709}}
  >{\raggedright\arraybackslash}p{(\columnwidth - 6\tabcolsep) * \real{0.2735}}@{}}
\toprule
\textbf{Indirect Pathway} & \textbf{Effect} & \textbf{Boot}\emph{SE} & \textbf{95\% CI} \\
 \midrule
X $\rightarrow$ M1 $\rightarrow$ Y ($a_1 b_1$) & 0.15 & 0.09 & {[}0.00, 0.34{]} \\
X $\rightarrow$ M2 $\rightarrow$ Y ($a_2 b_2$) & -0.60 & 0.18 & {[}-0.98, -0.28{]} \\
X $\rightarrow$ M1 $\rightarrow$ M2 $\rightarrow$ Y ($a_1 d_{21} b_2$) & -0.15 & 0.07 & {[}-0.32, -0.04{]} \\
Total indirect effect & -0.61 & 0.19 & {[}-1.01, -0.26{]} \\
Direct effect ($c'$) & 0.07 & 0.24 & {[}-0.39, 0.54{]} \\
\bottomrule
\end{tabular}
\end{table}

\noindent \textbf{Interpretive caveats.} This analysis rests on a single-item measure, was not preregistered, and captures only a reduction in one inference (revenue generation) rather than the full theorized shift in interpretation; increases in alternative inferences such as customer acquisition or community building were not measured. The results are mechanism-consistent but preliminary; confirmatory testing with multi-item measures is warranted.

\SIAnchorHeading{si:e}{E. Preliminary Field Study: OpenSea Floor-Price Analysis}

This section provides details of the marketplace analysis reported in the main text. As of April 2025, 77 brand-related collections were retrieved from OpenSea, of which 65 collections tied to the 32 focal brands were retained for analysis. Collections were classified as Free (\emph{N} = 37) or Paid (\emph{N} = 28) based on primary-market mint price. Descriptive statistics appear in Table E1.
% TODO(lineage): the full derivation from the 6,209-row raw export is being
% re-documented (99_documentation/DATA_AND_REPRODUCIBILITY_GAPS.md).
\setcounter{table}{0}
\renewcommand{\thetable}{E\arabic{table}}
\renewcommand{\theHtable}{E\arabic{table}}

\begin{table}[H]
\centering
\caption{Descriptive Statistics for Floor Price (Collection Level, ETH)}
\begin{tabular}{@{}
  >{\raggedright\arraybackslash}p{(\columnwidth - 4\tabcolsep) * \real{0.3333}}
  >{\raggedright\arraybackslash}p{(\columnwidth - 4\tabcolsep) * \real{0.3333}}
  >{\raggedright\arraybackslash}p{(\columnwidth - 4\tabcolsep) * \real{0.3333}}@{}}
\toprule
\textbf{Statistic} & \textbf{Free (\emph{N} = 37)} & \textbf{Paid (\emph{N} = 28)} \\
\midrule
Mean & 0.499 & 0.099 \\
Median & 0.014 & 0.017 \\
SD & 2.014 & 0.226 \\
Min & 0.001 & 0.002 \\
Max & 12.000 & 1.000 \\
\bottomrule
\end{tabular}
\end{table}

The higher mean for the Free group reflects outliers (e.g., AMC at 12.0 ETH); median floor prices were similar (Free = 0.014 ETH; Paid = 0.017 ETH). Test results appear in Table E2.

\begin{table}[H]
\centering
\caption{Tests Comparing Free vs. Paid Floor Prices}
\begin{tabular}{@{}
  >{\raggedright\arraybackslash}p{(\columnwidth - 4\tabcolsep) * \real{0.5556}}
  >{\raggedright\arraybackslash}p{(\columnwidth - 4\tabcolsep) * \real{0.2778}}
  >{\raggedright\arraybackslash}p{(\columnwidth - 4\tabcolsep) * \real{0.1667}}@{}}
\toprule
\textbf{Test / Model} & \textbf{Statistic} & \emph{p} \\
\midrule
Mann--Whitney \emph{U} & \emph{U} = 530.5 & .874 \\
Welch's \emph{t} (raw) & \emph{t} = 1.198 & .239 \\
Welch's \emph{t} (log) & \emph{t} = 0.556 & .580 \\
Mixed-effects model (log floor price; brand as random intercept) & $\beta$ = -0.22, \emph{SE} = 0.42, \emph{z} = -0.53 & .598 \\
Cohen's \emph{d} (log) & 0.137 & --- \\
\bottomrule
\end{tabular}
\end{table}

None of the analyses yielded a significant difference; we do not interpret these null results as evidence of equivalence. Within-brand comparison results are summarized in Table E3.

\begin{table}[H]
\centering
\caption{Within-Brand Analysis (Six Brands With Both Free and Paid Collections)}
\begin{tabular}{@{}
  >{\raggedright\arraybackslash}p{(\columnwidth - 4\tabcolsep) * \real{0.5000}}
  >{\raggedright\arraybackslash}p{(\columnwidth - 4\tabcolsep) * \real{0.2500}}
  >{\raggedright\arraybackslash}p{(\columnwidth - 4\tabcolsep) * \real{0.2500}}@{}}
\toprule
\textbf{Test} & \textbf{Statistic} & \emph{p} \\
\midrule
Paired \emph{t}-test & \emph{t} = 0.059 & .955 \\
Wilcoxon signed-rank test & \emph{W} = 10.0 & 1.000 \\
\bottomrule
\end{tabular}
\end{table}

The within-brand analysis included six brands that had both Free and Paid collections (36 collections; Free = 28, Paid = 8). In addition, 26 of the 28 Paid collections traded below their mint price.

\noindent \textbf{Limitations.} The unit of analysis is a cross-sectional collection-level dataset (\emph{N} = 65), which limits statistical power. Confounds that may affect secondary-market prices, such as supply volume, category, collection age, and market conditions, were not controlled.

\SIAnchorHeading{si:f}{F. Preliminary Field Study: Brand Classification (Luxury vs. Non-Luxury)}

This section reports the brand classification used in the X data analysis. The classification was fixed prior to the exploratory analysis and used to estimate the interaction with the Free/Paid classification.
\setcounter{table}{0}
\renewcommand{\thetable}{F\arabic{table}}
\renewcommand{\theHtable}{F\arabic{table}}

\begin{table}[H]
\centering
\caption{Brand Classification (32 Brands)}
\begin{tabular}{@{}
  >{\raggedright\arraybackslash}p{(\columnwidth - 2\tabcolsep) * \real{0.2137}}
  >{\raggedright\arraybackslash}p{(\columnwidth - 2\tabcolsep) * \real{0.7863}}@{}}
\toprule
\textbf{Category} & \textbf{Brands} \\
\midrule
Luxury (10) & Diesel, Gucci, KITH, Lacoste, Louis Moinet, McLaren, Mercedes-Benz, Porsche, Prada, Tiffany \& Co. \\
Non-Luxury (22) & AMC (The Walking Dead), Adidas, Austrian Post, Cameo, Havaianas, Hypebeast/Atmos, Hyundai, IMPS (The Smurfs), ITV Studios (Mr. Bean), ITV Studios (Thunderbirds), Jarritos, MEE6, Metro Trains Melbourne, Michelin, NFL, Nike/RTFKT, PUMA, Playboy, PokerGO, Square Enix, Sulake (Habbo Hotel), The Hundreds \\
\bottomrule
\end{tabular}
\end{table}

In the interaction analysis reported in the main text, Paid $\times$ Luxury was significant (\emph{$\beta$} = -0.41, \emph{p} \textless{} .001), indicating that the association between the Paid classification and lower VADER compound scores was stronger for the Luxury group.

\SIAnchorHeading{si:g}{G. Preliminary Field Study: X Post Analysis Procedures and Results}

This section supplements the X (formerly Twitter) analysis reported in the main text. The analytical sample comprised 22,841 English-language posts retrieved with NFT-collection-specific queries for collections of the same 32 brands (February 9, 2022 through April 28, 2025); the retrieval and cleaning pipeline is documented in Section I. Brand classification followed Section F (Luxury = 10, Non-Luxury = 22).

Each post was scored with VADER sentiment analysis (Hutto \& Gilbert, 2014), implemented in a manner consistent with prior online-review research (Costa et al., 2019). Estimation proceeded in two tracks: (1) a linear mixed model (LMM) with brand as a random intercept and calendar time as a covariate, and (2) a generalized estimating equations (GEE) model with negative posts (compound $\leq$ -0.05) as the outcome. The Free/Paid classification was coded as paid (0 = Free, 1 = Paid), and brand type was coded as luxury (0 = Non-Luxury, 1 = Luxury). GEE yields standard errors that are robust to within-brand correlation in the post-level data.

\noindent \textbf{Key results:}

Linear mixed model: The main effect of paid was significantly negative (\emph{$\beta$} = -0.03, \emph{SE} = 0.01, \emph{z} = -3.11, \emph{p} = .002). The main effect of luxury was significantly positive (\emph{$\beta$} = 0.32, \emph{SE} = 0.13, \emph{z} = 2.42, \emph{p} = .015). Calendar time was significant (\emph{$\beta$} = 0.0529, \emph{SE} = 0.0045, \emph{z} = 11.67, \emph{p} \textless{} .001).

Interaction: paid $\times$ luxury was significant (\emph{$\beta$} = -0.41, \emph{SE} = 0.12, \emph{z} = -3.48, \emph{p} \textless{} .001).

Simple slopes: The negative association with the Paid classification was larger for the Luxury group (\emph{$\beta$} = -0.44, \emph{SE} = 0.12, \emph{z} = -3.76, \emph{p} \textless{} .001) than for the Non-Luxury group (\emph{$\beta$} = -0.03, \emph{SE} = 0.01, \emph{z} = -3.11, \emph{p} = .002).

GEE (negative posts): The odds ratio for Paid was 1.24 (95\% CI {[}1.01, 1.51{]}, \emph{p} = .041).

Full coefficient estimates appear in Tables G1 and G2.
\setcounter{table}{0}
\renewcommand{\thetable}{G\arabic{table}}
\renewcommand{\theHtable}{G\arabic{table}}

\begin{table}[H]
\centering
\caption{X Post Analysis (Linear Mixed Model), Key Coefficients}
\begin{tabular}{@{}
  >{\raggedright\arraybackslash}p{(\columnwidth - 8\tabcolsep) * \real{0.3846}}
  >{\raggedright\arraybackslash}p{(\columnwidth - 8\tabcolsep) * \real{0.1282}}
  >{\raggedright\arraybackslash}p{(\columnwidth - 8\tabcolsep) * \real{0.1282}}
  >{\raggedright\arraybackslash}p{(\columnwidth - 8\tabcolsep) * \real{0.1282}}
  >{\raggedright\arraybackslash}p{(\columnwidth - 8\tabcolsep) * \real{0.2308}}@{}}
\toprule
\textbf{Term} & \emph{$\beta$} & \emph{SE} & \emph{z} & \emph{p} \\
\midrule
Intercept & 0.19 & 0.05 & 3.51 & \textless{} .001 \\
Paid (main effect) & -0.03 & 0.01 & -3.11 & .002 \\
Luxury (main effect) & 0.32 & 0.13 & 2.42 & .015 \\
Calendar time (covariate) & 0.0529 & 0.0045 & 11.67 & \textless{} .001 \\
Paid $\times$ Luxury (interaction) & -0.41 & 0.12 & -3.48 & \textless{} .001 \\
Simple slope (Non-Luxury) & -0.03 & 0.01 & -3.11 & .002 \\
Simple slope (Luxury) & -0.44 & 0.12 & -3.76 & \textless{} .001 \\
\bottomrule
\end{tabular}
\end{table}

\emph{Note.} Dependent variable: VADER compound score. Paid: 0 = Free, 1 = Paid. Luxury: 0 = Non-Luxury, 1 = Luxury. Brand (32 brands) entered as a random intercept (\emph{N} = 22,841 posts). Random-intercept variance ($\tau^2$) = 0.062; residual variance ($\sigma^2$) = 0.191; ICC = 0.245.

\begin{table}[H]
\centering
\caption{X Post Analysis (GEE With Negative Posts as Outcome), Key Results}
\begin{tabular}{@{}
  >{\raggedright\arraybackslash}p{(\columnwidth - 10\tabcolsep) * \real{0.3205}}
  >{\raggedright\arraybackslash}p{(\columnwidth - 10\tabcolsep) * \real{0.1175}}
  >{\raggedright\arraybackslash}p{(\columnwidth - 10\tabcolsep) * \real{0.1175}}
  >{\raggedright\arraybackslash}p{(\columnwidth - 10\tabcolsep) * \real{0.1175}}
  >{\raggedright\arraybackslash}p{(\columnwidth - 10\tabcolsep) * \real{0.2094}}
  >{\raggedright\arraybackslash}p{(\columnwidth - 10\tabcolsep) * \real{0.1175}}@{}}
\toprule
\textbf{Term} & \textbf{Coef} & \emph{SE} & \emph{z} & \textbf{OR {[}95\% CI{]}} & \emph{p} \\
\midrule
Intercept & -1.60 & 0.24 & -6.55 & --- & \textless{} .001 \\
Paid (Paid vs. Free) & 0.21 & 0.10 & 2.04 & 1.24 {[}1.01, 1.51{]} & .041 \\
\bottomrule
\end{tabular}
\end{table}

\emph{Note.} Dependent variable: VADER compound $\leq$ -0.05 (negative post = 1). GEE estimated with brand as the clustering variable under an exchangeable correlation structure (\emph{N} = 22,841, 32 brands). OR \textgreater{} 1 indicates that negative posts were more likely for the Paid classification.

\SIAnchorHeading{si:h}{H. Supplementary Study 4: Full Report (Favorable Displayed Information)}

\noindent \textbf{Purpose.} The five main-text experiments concern unfavorable or absent displayed outcome information. This supplementary study extends the program to favorable displayed information---a high buyer count and a rising price---after the same launch phase. Consistent with the main text, no formal hypothesis is stated; the study is reported for completeness and disclosure.

\noindent \textbf{Method.} A one-factor, two-level between-subjects experiment manipulated the launch condition (comparable price vs.\ free; wording in Section A-6), preregistered on AsPredicted (\#262807). Participants were adults aged 20 or older with an interest in NFTs, recruited through a crowdsourcing platform. Of 282 respondents, six who failed the attention check and three who met the preregistered ceiling criterion (a score of 10 at both \emph{T}0 and \emph{T}1) were removed, yielding \emph{N} = 273 (55.3\% male; median age bracket 40--44). The stimulus was the bee image with a small logo from Study 3b. Perceived Brand Luxury (three items; main-text Appendix) was measured at \emph{T}0 (baseline), \emph{T}1 (after the launch information), and \emph{T}2 (after the favorable displayed information: ``Order book (buyers for this NFT): 1,343 (very high)'' and a sharply rising price) using a continuous 0--10 slider suited to repeated within-subject measurement.

\noindent \textbf{Results.} A 2 (condition) $\times$ 3 (time point) mixed-design ANOVA was conducted; because the design involved two between-subjects groups, the error degrees of freedom for the within-subject factor and the interaction were (\emph{N} - 2) $\times$ (3 - 1) = 542. The main effect of time was significant (\emph{F}(2, 542) = 91.52, \emph{p} \textless{} .001, partial $\eta^2$ = .25). Bonferroni-adjusted pairwise comparisons showed that Perceived Brand Luxury at \emph{T}1 (\emph{M} = 5.53, \emph{SD} = 1.99) was lower than at \emph{T}0 (\emph{M} = 6.61, \emph{SD} = 1.67; \emph{t}(272) = 11.38, \emph{p} \textless{} .001, \emph{d} = 0.69), and that scores at \emph{T}2 (\emph{M} = 6.52, \emph{SD} = 2.04) rose from \emph{T}1 (\emph{t}(272) = -11.92, \emph{p} \textless{} .001, \emph{d} = -0.72); \emph{T}0 and \emph{T}2 did not differ significantly (\emph{t}(272) = 1.02, \emph{p} = .931, \emph{d} = 0.06). Neither the main effect of condition (\emph{F}(1, 271) = 0.03, \emph{p} = .856, partial $\eta^2$ \textless{} .001) nor the condition $\times$ time interaction (\emph{F}(2, 542) = 0.85, \emph{p} = .427, partial $\eta^2$ \textless{} .01) was significant. Condition-level means were: comparable \emph{T}0 = 6.66, \emph{T}1 = 5.46, \emph{T}2 = 6.49; free \emph{T}0 = 6.56, \emph{T}1 = 5.59, \emph{T}2 = 6.55; the conditions did not differ significantly at any time point (\emph{T}0: \emph{d} = 0.05, \emph{p} = .653; \emph{T}1: \emph{d} = -0.07, \emph{p} = .573; \emph{T}2: \emph{d} = -0.03, \emph{p} = .786), and the \emph{T}1-to-\emph{T}2 change was similar across conditions (comparable: \emph{M} = 1.03; free: \emph{M} = 0.96; \emph{d} = 0.05, \emph{p} = .683).

\begin{figure}[H]
\centering
\includegraphics[width=\linewidth]{figures/figure_6_study4.pdf}
\caption{Supplementary Study 4: Perceived Brand Luxury Across Time Points by Condition}
\end{figure}

\noindent \textbf{Discussion.} Perceived Brand Luxury declined after the launch information and rose after the favorable displayed information in both conditions. The design does not establish differential recovery between the comparable-price and free conditions, and the non-significant condition differences are not evidence of equivalence. The favorable information was a stylized cue (a simplified price display rather than an explicit separation of initial and current secondary-market prices), which bounds the generality of this pattern.

\SIAnchorHeading{si:i}{I. Additional Methodological Notes and Sensitivity Analyses}

\SIMajorHeading{I-1. Study 1: Equivalence Testing Details}

The main text reports that the non-significant difference between the free and comparable-price conditions on Loss of Brand Luxuriousness is not evidence of equivalence. The full two one-sided tests (TOST) results, computed on the attention-check-passed sample (\emph{N} = 670), are as follows. Free condition: \emph{n} = 166, \emph{M} = 3.906, \emph{SD} = 1.504; comparable condition: \emph{n} = 168, \emph{M} = 3.687, \emph{SD} = 1.595; mean difference (free -- comparable) = 0.219, Cohen's \emph{d} = 0.141, 90\% CI of the mean difference {[}-0.061, 0.499{]}; NHST two-sided \emph{p} = .198. At a smallest effect size of interest of \emph{d} = $\pm$0.30, the TOST upper-bound test was not significant (\emph{p} = .074), so equivalence was not established. Sensitivity analysis: equivalence was likewise not established at margins of $\pm$0.31 (\emph{p} = .062) or $\pm$0.32 (\emph{p} = .052), and was established at $\pm$0.33 (\emph{p} = .043), $\pm$0.35 (\emph{p} = .029), and $\pm$0.40 (\emph{p} = .009). The two conditions therefore cannot be distinguished at conventional significance, but effects up to roughly \emph{d} = 0.32 in either direction remain compatible with the data. For Behavioral Intention, the free and comparable conditions clearly differed (mean difference = 1.424, \emph{d} = 0.945, 90\% CI {[}1.152, 1.696{]}), so no equivalence question arises.

\SIMajorHeading{I-2. Scope of the Manipulation-Check Item in Studies 3b and 3c}

Studies 3b and 3c used a single check item, ``I think this NFT is the flagship product of the brand'' (adapted from Hubert et al., 2017). The item was written for the image contrast of Study 3b, where it showed a large difference between the central-GG-logo image and the bee image (\emph{d} = 1.57). In Study 3c the same item was reused for the source contrast, where the difference was smaller (\emph{d} = 0.39). The item establishes that the manipulated conditions differed on this judgment; it is a single item rather than a validated multi-item measure, and in Study 3c it is an indirect check of the source manipulation, because the manipulated variable is where the launch and outcome were presented rather than the offering's product status. Multi-item and manipulation-specific checks are a priority for future data collection.

\SIMajorHeading{I-3. Implemented Free and Complimentary Conditions Across Studies}

The free or complimentary condition was implemented differently across studies (main-text Table 2): plain free distribution in Study 1; a complimentary gift with a fashion magazine in Studies 2 and 3a (chosen because an outright giveaway by a luxury brand lacked scenario plausibility for a physical product); a complimentary gift with a fan book in Study 3b; and plain free distribution in the text-only Study 3c scenarios. One preregistration deviation follows from this: the Study 3c preregistration described the free condition as a fan-book bonus, whereas the implemented scenario presented plain free distribution. Two consequences for interpretation are drawn consistently in the main text. First, each study's result is interpreted at the level of its implemented contrast, and the implementations are never pooled as a single treatment. Second, the recurrence of the first-stage pattern of H3a across these different implementations (Study 3a; the bee image in Study 3b; the news article in Study 3c) is read as conceptual replication of the direction of the effect, not as replication of a single quantitative effect.

\SIMajorHeading{I-4. X Data: Retrieval Pipeline and Audit Status}

Posts were retrieved with NFT-collection-specific queries combining collection and project names and hashtags with NFT/Web3 disambiguation terms; each post was mapped to a collection slug and inherited that collection's Free/Paid classification. Retrieval returned 33,530 posts, of which 27,467 were English-language. Sequential cleaning steps and their effects on the sample were as follows.

\begin{table}[H]
\centering
\begin{singlespace}
\caption*{Cleaning Pipeline for the X Post Sample}
\begin{tabular}{@{}
  >{\raggedright\arraybackslash}p{(\columnwidth - 4\tabcolsep) * \real{0.60}}
  >{\raggedleft\arraybackslash}p{(\columnwidth - 4\tabcolsep) * \real{0.18}}
  >{\raggedleft\arraybackslash}p{(\columnwidth - 4\tabcolsep) * \real{0.18}}@{}}
\toprule
\textbf{Step} & \textbf{Removed} & \textbf{Remaining} \\
\midrule
Start (English-language posts) & --- & 27,467 \\
Duplicate post IDs (same post retrieved for multiple collections) & 679 & 26,788 \\
Accounts with ``bot'' in the username (predominantly mechanical sale/listing notifications) & 3,331 & 23,457 \\
Duplicate post texts & 577 & 22,880 \\
URL-only posts & 4 & 22,876 \\
Posts under 20 characters & 29 & 22,847 \\
Brands with fewer than 10 posts & 6 & 22,841 \\
\bottomrule
\end{tabular}
\end{singlespace}
\end{table}

The final sample comprised 22,841 posts across 32 brands (Free 10,381; Paid 12,460). Three audit points remain open and are disclosed in the main text. First, official brand and project accounts were not separately excluded; an account-level audit is planned. Second, query precision and duplicate assignment across collections have not been fully audited beyond the duplicate-ID removal above. Third, the models cluster posts by brand, whereas the Free/Paid classification varies at the collection level; collection-level clustering is a target of ongoing reanalysis. In addition, only three luxury brands contributed Free-classified collections (and only one brand contributed both Free and Paid collections), so the Luxury interaction reported in the main text is estimated primarily from between-brand comparisons and is presented as exploratory heterogeneity.
% TODO(lineage): counts and audit status to be finalized against the
% reproducibility audit (99_documentation/DATA_AND_REPRODUCIBILITY_GAPS.md).

\newpage
\noindent \textbf{References}
\begin{referenceshang}

Costa, H., Gil, R., \& Rita, P. (2019). Unfolding the characteristics of incentivized online reviews. \emph{Journal of Retailing and Consumer Services, 52}, 101836. \url{https://doi.org/10.1016/j.jretconser.2018.12.006}

DerSimonian, R., \& Laird, N. (1986). Meta-analysis in clinical trials. \emph{Controlled Clinical Trials, 7}(3), 177--188. \url{https://doi.org/10.1016/0197-2456(86)90046-2}

Hayes, A. F. (2022). \emph{Introduction to mediation, moderation, and conditional process analysis: A regression-based approach} (3rd ed.). Guilford Press.

Hubert, M., Florack, A., Gattringer, R., Eberhardt, T., Enkel, E., \& Kenning, P. (2017). Flag up!\,--- Flagship products as important drivers of perceived brand innovativeness. \emph{Journal of Business Research, 71}, 154--163. \url{https://doi.org/10.1016/j.jbusres.2016.09.001}

Hutto, C. J., \& Gilbert, E. (2014). VADER: A parsimonious rule-based model for sentiment analysis of social media text. \emph{Proceedings of the International AAAI Conference on Web and Social Media, 8}(1), 216--225. \url{https://doi.org/10.1609/icwsm.v8i1.14550}

McShane, B. B., \& B\"ockenholt, U. (2025). Single paper meta-analysis is unavoidable. \emph{Journal of Consumer Psychology, 35}(4), 663--685. \url{https://doi.org/10.1002/jcpy.1462}
\end{referenceshang}
\end{document}
